% sections/ndss_08_related.tex  — NDSS 2027

\paragraph{Prompt injection.}
Greshake \etal~\cite{greshake2023indirect} introduce indirect prompt injection
and demonstrate that LLM agents can be manipulated through
attacker-controlled content in retrieved documents.
Our work complements this: we show that when dangerous sink capabilities
are undisclosed, even a user who is aware of prompt injection
cannot protect themselves---they cannot know which installed skills
can be weaponized as sinks.
Recent work on multi-agent injection~\cite{wu2024autogen}
and tool-call hijacking~\cite{debenedetti2024agentdojo}
further establishes the attack surface; we provide the first
marketplace-scale measurement of the exposed capability surface.

Table~\ref{tab:injection_comparison} formally contrasts our threat model
with classical prompt injection and plugin-level attacks:

\begin{table}[h]
\centering
\caption{Source-sink composition attack vs.\ related threat models.}
\label{tab:injection_comparison}
\scriptsize
\setlength{\tabcolsep}{3pt}
\begin{tabular}{lllll}
\toprule
 & \textbf{Classic PI} & \textbf{Plugin atk} & \textbf{This work} \\
\midrule
Attacker controls & Input & Plugin code & Source skill output \\
Exploited gap & Model trust & Code execution & Capability omission \\
Requires runtime & Yes & Yes & No (install-time) \\
User can prevent? & Careful inputs & Vetted plugins & Not without disclosure \\
Sink needed? & No & No & Yes (IC-1 skill) \\
\bottomrule
\end{tabular}
\end{table}

\paragraph{Mobile app permission analysis.}
Extensive work on Android permissions~\cite{felt2011android,au2012pscout}
and privacy policy consistency~\cite{andow2019policylint,bui2021consistency}
motivated our methodology.
Liao \etal~\cite{liao2020measuring} measure over-privilege in Alexa skills
using a difference matrix approach; we adapt this to LLM agent marketplaces
where the permission model is replaced by natural-language descriptions.
Key differences: skills have richer implementation code (not just intents),
can execute arbitrary shell commands, and operate with full OS-level authority.

\paragraph{Software supply chain attacks.}
Ohm \etal~\cite{ohm2020backstabber} and Ladisa \etal~\cite{ladisa2023taxonomy}
taxonomy malicious packages in npm/PyPI.
LLM skill supply chain attacks differ in that the malicious payload
delivery mechanism is NLP-mediated: the LLM decides whether to call
a tool based on context, not static control flow.
This makes both detection (static analysis alone is insufficient)
and exploitation (prompt injection enables remote triggering) qualitatively different.

\paragraph{LLM agent security.}
Recent work studies jailbreaking~\cite{wei2023jailbroken},
tool misuse~\cite{ruan2024identifying},
and multi-agent trust~\cite{ju2024flooding}.
\sysname focuses on the marketplace supply chain rather than
the agent runtime---the attack surface created before any
user interaction, at skill install time.
